\subsection{Adapter: Windowed Capacity and Layer Control}
\label{sec:kv-runtime}
\label{sec:capacity-allocation}

Adapter (4) turns the KV Master's policy into physical copies and
recovery execution. Control windows last 5\,s. All agents' requests for
the same deployed model share one target resident depth $L_t$, unchanged
throughout window $t$, including requests arriving during that window.
Each target DP controls its own total Peer-HBM budget $B_{p,t}$ and
applies this common depth to admitted Peer requests. DRAM-only and
recomputation tiers do not require Peer residency. Adapters report actual
residency and commit status; recovery uses committed objects, since
copying may not yet have reached the target depth.

% Target-local capacity with a common window-level depth, part of section 3.6.
\paragraph{Control Peer capacity each window.}
For a target DP $p$ in window $t$, its Adapter grants
\begin{equation}
 \begin{aligned}
 B_{p,t}=\bigl[\min\bigl(&B_{\mathrm{cap},p,t},\\
 &\mathrm{pool}_{p,t}-\mathrm{watermark}_{p,t}
              -\mathrm{headroom}_{p,t}\bigr)\bigr]^+,
 \end{aligned}
 \label{eq:peer-budget}
\end{equation}
where $B_{\mathrm{cap},p,t}=\tfrac12\mathrm{pool}_{p,t}$ limits Peer
backups to half of the managed KV pool, after model weights and fixed
runtime allocations. Foreground demand and headroom can reduce the
grant further. The pool includes bytes already occupied by Peer replicas;
$B_{p,t}$ is a total target, not additional free space. Foreground demand,
admission headroom, and transfer/compute estimates are refreshed each window.

Let $A_r(L_t)$ contain the KV Master's current hot sets in layers
$1\ldots L_t$ and their required selector state. The Adapter processes
$Q_p$ in the Master's priority order. Starting with $P_s=\varnothing$,
it admits a request's complete footprint only if
\begin{equation}
 \mathrm{bytes}\bigl(P_s\cup A_r(L_t)\bigr)\le B_{p,t},
 \label{eq:resident-prefix}
\end{equation}
then adds $A_r(L_t)$ to $P_s$. Shared physical objects count once.
If it does not fit, the request retains its DRAM path and the Adapter
considers the next candidate. All admitted requests use the same $L_t$;
their hot-block counts and byte footprints can differ. This greedy
admission has no global optimality claim and does not search over
request-specific depths.

Within a window, $L_t$ is fixed, while sealing can enlarge a request's
hot sets. Growth reserves additional bytes before copying. If capacity
is insufficient, admission is reconsidered in the Master's priority
order; emergency foreground pressure withdraws whole Peer footprints
from the lowest-priority requests, preserving DRAM recovery. Neither
case changes $L_t$. Copying, committed, and reclaiming bytes remain
physically charged until release completes. New copies wait for physical
headroom, readers retain their pins, and growth transfers only newly
required objects. Target changes do not instantly free addresses or
make new replicas readable.


\paragraph{Inform the common depth with a reference profile.}
The Adapter preloads hot blocks in layer order, skipping Peer-resident
objects. Selector state must arrive before exact demand can be computed;
newly revealed misses take priority over speculative preloads. Cold KV
is filled in the background. Let $h$ be the configured preload fraction and
$\widetilde u$ the fraction of historical bytes in the union of preload
and demand-fetch sets, calibrated from replay traces. For $N$ homogeneous
layers, per-layer replay compute $C$ and full-layer transfer
$D=M_rs/\beta$ ($M_r$ historical tokens, $s$ KV bytes/token/layer, and
$\beta$ effective bandwidth),
\begin{align}
 W_{\mathrm{floor}}&=N(\widetilde u-h)D,\label{eq:demand-floor}\\
 W_0&=hD+(N-1)[hD-C]^++W_{\mathrm{floor}}.\label{eq:exposed-wait}
\end{align}
Here $[x]^+=\max(0,x)$. The first two terms are startup preload and
backlog that computation cannot hide. The last is demand wait under a
fixed hot set, causal demand release, and ideal preemptible copy stream;
it is not a lower bound on other hot sets or loading policies.

For a configured representative history and replay tail, the window's
transfer/compute estimates give a reference preload-byte target
\begin{equation}
 F^\star_{\mathrm{ref},t}=\beta_t
 \bigl(W_{0,\mathrm{ref},t}-W_{\mathrm{floor},\mathrm{ref},t}\bigr).
 \label{eq:peer-byte-target}
\end{equation}
The smallest $k$ whose reference hot sets satisfy
$\sum_{l=1}^{k}K_{\mathrm{ref},l}b\ge F^\star_{\mathrm{ref},t}$
provides the common $L_t$ at the window boundary; $b$ is KV bytes per
layer-block. The depth is zero when no preload is useful and capped at
$N$. The reference profile is a shared calibration choice, not a
separate optimization for each request. Selector-state bytes are
additionally charged in $A_r(L_t)$. When $hD\le C$, one layer's hot set
suffices for this homogeneous reference. Different request histories
can still have different residual waits: a common depth does not promise
to reach every request's floor. Capacity controls request admission
through Equation~\ref{eq:resident-prefix}, without shortening individual
requests' target depths inside the window.

\begin{figure}[t]
  \centering
  \includegraphics[width=\columnwidth]{figures/design-single/design-layer-pipeline.pdf}
  \caption{Modeled recovery pipelines: MiniCPM4-8B, 32 layers, 32K
    history, 64-token known tail; $C=2.87$\,ms/layer (measured) and
    $D=5.37$\,ms at an assumed 6.25\,GB/s. In (d), 147\,MiB of Peer
    residency removes preload wait. Totals end at tail-KV reconstruction.
    The homogeneous, preemptible-copy model excludes detection,
    queuing, and interference.}
  \label{fig:design-pipeline}
  \Description{Four timelines compare token-block transfer, full-layer
    pipelining, hot preloads with demand fetch, and Peer-resident hot sets.}
\end{figure}

Figure~\ref{fig:design-pipeline}(a,b) reduces exposed wait from 172 to
83\,ms through layer pipelining; hot preloads in (c) reduce it to 44\,ms;
seven resident hot sets in (d) leave 20\,ms of demand wait. This is one
calibrated profile illustrating the useful prefix, not a fixed layer
count across windows or measured fault recovery.

\paragraph{Publish only complete replicas.}
Adapters invoke Connectors to allocate, transfer, and acknowledge
objects using the assigned targets and the current $D_s,P_s$.
Peer and fault-independent DRAM paths commit separately. Publication
follows transfer completion and model/layout/shard validation.
Selector state must be ready before selection. A Peer object that is
the only surviving copy cannot be discarded before its DRAM copy
commits. Low-priority copy streams are rate-limited (default 20\% of
link capacity), also accounting for the target's HBM bandwidth.

\paragraph{Preserve attention inputs during recovery.}
For each retry, preloads overlap earlier-layer computation; newly
revealed misses outrank queued preloads. The runtime's DMA granularity
may limit ideal preemption. A layer executes only when \emph{all} its
required attention inputs are ready; later cold blocks can arrive in
the background or be fetched by subsequent decode. Tier R rebuilds
the history. AegisKV changes placement and loading order, not selector
outputs or known replay tokens. With compatible weights, layout, and
operators it reconstructs the same logical KV state without an added
attention approximation, but does not guarantee identical future
stochastic samples across different floating-point schedules.
